$\phi$ is an algebra map because $\varepsilon_+$ is a central idempotent of $A$: $\phi(a)\phi(b) = \varepsilon_+ a\, \varepsilon_+ b = \varepsilon_+^2 ab = \phi(ab)$, and $\phi(1) = \varepsilon_+$. Injectivity: if $a$ is even and $\varepsilon_+ a = 0$, then $a = (\varepsilon_+ + \varepsilon_-)a = \varepsilon_- a$, i.e.\ $2a = a - \omega a$, i.e.\ $\omega a = -a$. But $\omega$ has odd grade ($n$ odd) while $a$ is even, so $\omega a$ is odd and $-a$ is even; a common value is possible only if both vanish, so $a = 0$. Surjectivity follows by dimension: $\dim \Cl^0 = 2^{n-1} = \tfrac12 \dim A = \dim \varepsilon_+ A$ (the first equality is \claimref{std-algebra.z2-grading}), the last equality because $x \mapsto (\varepsilon_+ x, \varepsilon_- x)$ splits $A$ into two summands exchanged (linearly) by the grade involution, which maps $\varepsilon_+ x \mapsto \varepsilon_- \hat{x}$ and is bijective. For the second statement: $\varepsilon_+ B_\gamma \varepsilon_+ = \varepsilon_+ A \varepsilon_+ + \varepsilon_+ A \nu \varepsilon_+$ by Lemma \claimref{transport.decomposition-commutation}; the first term is $\varepsilon_+ A$ (centrality in $A$), and the second vanishes: $\nu \varepsilon_+ = \varepsilon_- \nu$ (Proposition \claimref{transport.ascent-renders-grade-involution-inner}), so $\varepsilon_+ a \nu \varepsilon_+ = \varepsilon_+ a \varepsilon_- \nu = (\varepsilon_+ \varepsilon_-) a \nu = 0$.
